\begin{onehalfspace}

Questo lavoro ha affrontato una domanda operativa: quanta informazione
dell'output rumoroso di Shor può essere recuperata a valle e quando diventa
necessario correggere l'errore durante il calcolo? La risposta delimita i due
regimi. Su $N=15$ una ricerca multi-candidato sfrutta meglio l'informazione già
misurata, mentre un classificatore applicato allo stesso output non aggiunge
valore. Oltre questo margine occorre intervenire sulle sindromi e sulla
correzione d'errore.

% -----------------------------------------------
\section{Le Domande di Ricerca}
\label{sec:risposte_dr}

\paragraph{DR1 --- Come cambia Shor al variare dei canali di rumore?}
Nel modello uniforme per $N=15$, aumentando $\lambda_{2q}$ la massa sui quattro
picchi QPE scende da $0.804$ a $0.042$, mentre il proxy di nessun evento Pauli
sulle 224 CX decade da $0.811$ a $2.65\times10^{-10}$: quest'ultimo non predice
quindi il successo algoritmico. Gli sweep di $\lambda_{1q}$, $T_1/T_2$, readout
e shot non ordinano causalmente i canali con 30 repliche; TOP-4 resta invece a
una iterazione, salvo il valore $1.07$ a $\lambda_{2q}=0.1$. Risultati e
diagnostiche complete sono nelle Sezioni~\ref{sec:parametri_mitigazione}
e~\ref{app:parametrica}.

\paragraph{DR2 --- Come si estrae una sindrome senza misurare lo stato logico?}
Nel codice a ripetizione si misurano le parità $Z_1Z_2$ e $Z_2Z_3$, non i qubit
logici: la sindrome localizza il bit alterato senza rivelarne l'ampiezza. La
curva Monte Carlo coincide con $p_L=3p^2-2p^3$ entro l'incertezza, ma protegge
un solo tipo di errore per volta (Sezione~\ref{sec:repetition}).

\paragraph{DR3 --- Quali vantaggi e limiti offre Steane $[[7,1,3]]$?}
La struttura CSS corregge un errore Pauli arbitrario su un qubit e la
corrispondenza sindrome--errore è verificata nei 21 casi. La pendenza log--log
$1.94$ conferma la soppressione quadratica attesa; distanza fissa ed estrazione
non fault-tolerant ne impediscono però una lettura scalabile
(Sezione~\ref{sec:steane}).

\paragraph{DR4 --- Il surface code mostra un comportamento di soglia?}
Sì, entro il memory experiment e il modello circuit-level adottati: le curve
per $d=3,5,7,9$ si incrociano intorno allo $0.9\%$ in base Z e allo $0.7\%$ in
base X, con soppressione sotto soglia e amplificazione sopra soglia. Il relativo
$p$ non coincide con $\lambda_{2q}$ della campagna Shor, perciò il confronto è
solo qualitativo (Sezione~\ref{sec:surface_risultati}).

\paragraph{DR5 --- Quale errore logico è compatibile con Shor?}
Nel modello fenomenologico, a $p_L=0$ il successo per esecuzione resta circa
$0.75$ per la quota intrinsecamente non utile degli esiti QPE; le ripetizioni ne
aumentano la probabilità cumulativa. La curva è però una sensibilità, non una
previsione end-to-end: il logical error rate di una memoria non può essere
assegnato a ogni gate di Shor senza operazioni fault-tolerant e canali residui
espliciti (Sezione~\ref{sec:risultati_logico}).

\paragraph{DR6 --- Quando un decoder appreso può superare quello analitico?}
Il vantaggio compare quando il rumore contiene correlazioni assenti dal grafo
del matching. Il decoder ibrido riduce l'errore rispetto a \ac{MWPM} a distanza
3, ma il margine si restringe e non è significativo a distanza 7 per la rete
densa adottata. Il risultato individua quindi una nicchia residuale, non una
sostituzione generale dei decoder analitici (Sezione~\ref{sec:decodifica_appresa}
e Appendice~\ref{app:decodifica}).

% -----------------------------------------------
\section{Verifica delle Ipotesi}
\label{sec:verifica_ipotesi}

\begin{table}[H]
\centering
\footnotesize
\setlength{\tabcolsep}{4pt}
\caption[Esito delle ipotesi]{Esito dopo la rigenerazione schema~2.}
\label{tab:esito_ipotesi}
\begin{tabular}{p{3.8cm}p{2.2cm}p{7.2cm}}
\toprule
\textbf{Ipotesi} & \textbf{Esito} & \textbf{Evidenza} \\
\midrule
Il classificatore riduce le iterazioni
& smentita
& M2 non migliora M1; TOP-4 supera M2 in UC1 e UC2
($p_{\mathrm{Holm}}=0.0033/0.0096$). \\
\addlinespace
F1 $>0.85$ e AUC $>0.90$
& confermata, ma insufficiente
& SVM: F1 $0.919/0.923$ e AUC $0.965/0.958$; le metriche descrivono TOP-1,
non l'utilità operativa. \\
\addlinespace
Il vantaggio si mantiene crescendo $N$
& non verificata
& $N=21/35$ sono validati idealmente ma esclusi dal rumore per costo; va
controllato anche l'ordine $r$. \\
\addlinespace
Soppressione quadratica a distanza 3
& confermata nel modello
& Ripetizione coerente con $3p^2-2p^3$; Steane con pendenza $1.94$. \\
\addlinespace
Soglia del surface code
& confermata nel modello
& Incrocio delle curve $d=3,5,7,9$ e inversione del beneficio sopra soglia. \\
\addlinespace
Trasferimento diretto di $p_L$ a Shor
& non dimostrato
& La curva è fenomenologica; manca una mappatura per operazione
fault-tolerant. \\
\bottomrule
\end{tabular}
\end{table}

L'audit chiarisce perché buone metriche non bastano: TOP-1 genera classi
addestrabili, mentre TOP-16 rende positivi tutti i $2\,000$ campioni di entrambi
gli scenari. Con 63 esiti utili su 256, la probabilità combinatoria di non
includerne alcuno fra sedici celle è appena
$\binom{193}{16}/\binom{256}{16}=0.9275\%$. La costruzione delle etichette e i
relativi controlli sono documentati nell'Appendice~\ref{sec:audit_etichette}.

% -----------------------------------------------
\section{Contributi del Lavoro}

\textbf{Demo ideale e rumorosa separate.}
Il percorso ideale espone circuito, stato e QPE; quello rumoroso rende
selezionabili depolarizzazione 1Q/2Q, $T_1/T_2$, readout ed errore coerente.
$N=21/35$ restano nel perimetro ideale e $N=15$ in quello rumoroso.

\textbf{Ablazione appaiata e riproducibile.}
M1, TOP-4 e M2 ricevono lo stesso istogramma. TOP-4 porta $\bar M$ da $1.37$ a
$1.00$ nello scenario moderato e da $1.50$ a $1.00$ nello stress, mentre il
classificatore peggiora la propria ablazione; seed, tie-break, sentinelle,
revisioni e manifest sono conservati negli artefatti.

\textbf{Validazione funzionale dell'aritmetica.}
Il moltiplicatore $\times7\bmod15$ è coperto da truth table; il percorso
Beauregard per $N=21/35$ da proprietà esaustive e confronti QPE. Il controllo
mostra perché picchi o ordine, isolatamente, non validano l'operatore.

\textbf{Progressione QEC con limiti espliciti.}
Ripetizione, Steane, surface code e decodifica appresa sono collegati da
metriche dichiarate senza rendere commensurabili modelli di rumore diversi;
Shor logico chiude la progressione come analisi fenomenologica.

% -----------------------------------------------
\section{Limiti}
\label{sec:limiti_finali}

Nessun risultato deriva da hardware quantistico. È una scelta concordata: la
disponibilità di un dispositivo è sopraggiunta dopo la chiusura delle domande
sui modelli di rumore e sulle snapshot di calibrazione. La validazione resta
quindi simulativa; il solo riscontro su macchina reale è quello discusso per
citazione nella Sezione~\ref{subsec:discussione_logico}.

La campagna Shor usa rumore uniforme, indipendente e stazionario, senza
topologia, routing, crosstalk, leakage o deriva. $N=15$ è didattico, TOP-K ha
una verifica classica permissiva, gli sweep impiegano 30 repliche e i test non
sono corretti globalmente fra otto famiglie; $N=21/35$ non sono simulati sotto
rumore. Nella QEC, Steane non ha estrazione fault-tolerant, memoria e proxy per
gate non condividono lo stesso $p_L$, e il decoder denso non mantiene il
vantaggio alle maggiori distanze. Le stime crittografiche richiedono dunque
architetture fault-tolerant complete, non estrapolazioni dalle curve di $N=15$.

% -----------------------------------------------
\section{Il Criterio che Unifica il Lavoro}

Un complemento semplice è utile quando sfrutta informazione trascurata dal
metodo base senza alterarne la generazione: TOP-4 verifica candidati già
nell'istogramma e il decoder residuale interviene dove il matching ha una
debolezza rappresentazionale. La complessità non crea valore se replica una
regola semplice, usa etichette degeneri o costa più del beneficio. Ne segue il
confine fra mitigazione e correzione: TOP-K recupera un picco già misurato, ma
non ricostruisce informazione distrutta; la scalabilità richiede operazioni
logiche fault-tolerant e l'apprendimento è utile solo se dispone di struttura
assente nel decoder analitico.

L'AQFT applica lo stesso criterio a monte. Troncare le rotazioni interne agli
addizionatori dimezza le porte a due qubit ($49\,661\rightarrow23\,178$), ma su
$N=21$ riduce la resa ideale di circa $6.5$ punti; il vantaggio massimo, circa
tre punti, richiederebbe porte quasi duecento volte più accurate, mentre il
circuito resta circa 150 volte più lungo della coerenza disponibile. Nella
stima di fase, molto più corta, la variante troncata vince invece tutti i nove
confronti e raggiunge $+0.126$ a dieci qubit. Il grado ottimo resta due o tre,
il risparmio cresce con il registro e le fasi esattamente rappresentabili
beneficiano di più: ciò spiega la differenza fra $r=4$ per $N=15$ e $r=6$ per
$N=21$. Conteggi compilati e calibrazione permettono quindi di scegliere il
troncamento prima dell'esecuzione. Dati, limite teorico e perimetro di
generalizzazione sono nella Sezione~\ref{sec:qft_topologia}.

Questa riduzione delle porte è complementare, ma non automaticamente additiva,
alla riduzione dei qubit di lavoro ottenuta con aritmetica approssimata
\cite{chevignard2025,gidney2025} e alla riduzione del registro di fase mediante
blocchi indipendenti~\cite{shukla2026}. Le tre leve agiscono su porte, qubit di
lavoro e qubit di conteggio; studiarne congiuntamente costi e profondità è la
domanda aperta più naturale.

\end{onehalfspace}
